\documentclass[12pt]{article}

\usepackage[a4paper, total={6in, 8in}]{geometry}
\usepackage{textcomp}

\begin{document}
\setlength{\parindent}{0pt}

\def \t {\textrightarrow}
\def \v {\vspace{1em}}
\def \bi {\begin{itemize}}
\def \ei {\end{itemize}}
\def \s[#1] {\section*{#1}}
\def \ss[#1] {\subsection*{#1}}
\def \sss[#1] {\subsubsection*{#1}}

\s[Oscar Wilde]
He is considered part of an international literary movement known as "aestheticism" \t it originated in France, in the mid 19th century, and then spread all around Europe \\
In Britain Walter Pater brought it and popularized it \t he was an intellectual and was Wilde's professor at Oxford university \\
However, Wilde is considered the most relevant member of this movement in Britain \\
It is based on the principle of "Art for art's sake", according to which literature (and art in general) has got no practial purpose, and should not be produced to convey a moral message or to teach something \\
Art is a mere expression of beauty \t art should have no moral and no political commitment \t and realism does not have to be a feature \\
Aestheticism is not a victorian movement \t victorians focused on moralism, utilitarianism and realism

\v

Oscar Wilde is also a member of the "Decadence", which is more of a cultural atmosphere than a movement (as something that developed in a non planned way) \\
Decadence has two meanings:
\bi
  \item because of the scandal decadent writers were creating with their literary production and lifestyle \t it happens with Oscar Wilde
  \item decadent writes were giving voice to a sense of ending for the age they were living in and for its middle-class hypocritical values \t the victorian society was going to finish
\ei

\ss[The happy prince]
It was not intended for publication, it was intended as a bet time story for his children \\
The city is unknown, but there is a port \t and there is a colossal statue of the happy prince \\
Questions:
\begin{enumerate}
  \item it was beautufil and full of gold and diamonds
  \item people appriciated it for its beauty and for the moral message it was conveying \t one has to enjoy life
  \item a little bird landed on the statue and he was flying alone, headed to Egypt
  \item he had a privileged life and did not know about suffering
  \item he was crying because he had to see the suffering of the people of the city
  \item to become his messanger with the aim to bring parts of him to poor people
  \item she was a seamstress and her son was ill, and in order to take care of him he could not work properly
  \item he helps a writer, who cannot concentrate because of the cold (he is poor)
  \item helps then a little girl whose eggs she was selling got broken and father would have beaten her
  \item in heaven, the spirits of the bird and the prince meet in the presence of god
\end{enumerate}
Moral messages are generosity, happiness comes from sharing, appearence does not matter, + religious messages \t so it is not really realted to the principle of aestheticism \\
However there is a refernce to beauty and art, which is typical of aestheticism

\ss[Picture of Dorian Gray]
The first edition was published in the year 1890 on an American magazine (not in the UK) \t it was published as a volume in the UK in 1891, together with a preface which is extremly relevant, and is considered the manifesto of british aestheticism \\
While the preface of the "Lyrical ballads" was the manifesto of british romanticism \\
Plot on page 126 \t ASK NOTES ABOUT PLOT
\end{document}
